There are several different methods which have been used for analyzing transients in \glspl{MSR}.
These methods include multi-group neutron diffusion, \gls{PRK} schemes, and quasi-statics \cite{wooten2018review}.
These methods are generally coupled, either tightly or loosely, to a thermal-hydraulics solver.
In general, although these methods have varying accuracy and costs, none of them include \gls{MSR} specific phenomena when generating the \gls{DNP} group parameters.

\subsubsection{Multi-Group Diffusion}
\label{sec:mgd-kinetics}

Multi-group diffusion kinetics methods are among the most accurate, and computationally expensive, approaches applied to transient problems, surpassed only by transport methods.
Hurwitz and Ehrlich provide a standard formulation of multi-group diffusion written in 1953  \cite{hurwitz1953multigroup}.

Examples of software that use fluid-flow-velocity field coupled multi-group include GenFOAM and Griffin/Pronghorn coupling using \gls{MOOSE} \cite{fiorina2015gen, tano2023coupled}.
Other examples can be found in the review by Wooten at al. \cite{kvrepel2007dyn3d, fiorina2014modelling, kophazi2009development} and elsewhere in the literature \cite{zhang2009analysis, yamamoto2006transient, nicolino2008coupled, park_advancements_2025}.
%Even more studies which use the multi-group diffusion method can be found in the literature \cite{zhang2009analysis, yamamoto2006transient, krepel2007dyn3d, nicolino2008coupled, park_advancements_2025}.
These studies typically use Picard iteration with multi-group neutron diffusion coupled with thermal-hydraulics to accurately capture the multi-physics effects of the reactor.
However, this accuracy also comes with a higher computational cost than simpler methods.
%The general form of the coupling can be seen in Equations \eqref{eq:mg-n-diff} and \eqref{eq:mg-n-diff-Ci} \cite{wooten2018review}.
The general form of the coupling is given as \cite{wooten2018review}:

\begin{equation}
    \begin{split}
    \frac{1}{\bar{\nu}_g} \frac{\partial \phi_g}{\partial t} = & \; \nabla \cdot D_g \nabla \phi_g + \sum_{g' \neq g}^G \Sigma_{s, g' \rightarrow g} \phi_{g'} \\
    & + (1-\beta) \chi_{p, g} \sum_{g'=1}^{G} (\nu \Sigma_f)_{g'} \phi_{g'}\\
    & \sum_{i=1}^I \chi_{d, g, i} \lambda_i C_i - \Sigma_{a, g} \phi_g - \sum_{g' \neq g}^G \Sigma_{s, g \rightarrow g'} \phi_{g},
    \label{eq:mg-n-diff}
    \end{split}
\end{equation}
and,
\begin{equation}
    \frac{\partial C_i}{\partial t} = \beta_i \sum_{g=1}^G (\bar{\nu}_g \Sigma_f)_g \phi_g - \lambda_i C_i - \nabla \cdot (C_i \vec{u}) + \nabla \cdot \frac{\nu_T}{Sc_T} \nabla C_i,
    \label{eq:mg-n-diff-Ci}
\end{equation}
where
\begin{align}
\bar{\nu}_g &= \mbox{ the average number of neutrons produced by fission of group \textit{g} [-]} \nonumber\\
\phi_g &= \mbox{ the scalar neutron flux of group \textit{g} [$n \cdot cm^{-2} \cdot s^{-1}$]} \nonumber\\
D_g &=\mbox{ the group \textit{g} diffusion term [$cm$]}\nonumber\\
G &= \mbox{ the total number of energy groups [-]} \nonumber\\
\Sigma_{s, g' \rightarrow g} &= \mbox{the macroscopic scattering cross section from group \textit{g'} to group \textit{g} [$cm^{-1}$]} \nonumber\\
\chi_{p, g} &= \mbox{the number of prompt neutrons born in group \textit{g} [-]} \nonumber\\
\chi_{d, g, i} &= \mbox{the number of delayed neutrons from precursor group \textit{i} born in group \textit{g} [-]} \nonumber\\
C_i &= \mbox{the concentration of \glspl{DNP} in group \textit{i} [$atoms \cdot cm^{-3}$]} \nonumber\\
\vec{u} &= \mbox{the velocity profile [$cm \cdot s^{-1}$]} \nonumber\\
\Sigma_{a, g} &= \mbox{the macroscopic absorption cross section for group \textit{g} [$cm^{-1}]$} \nonumber\\
\nu_T &= \mbox{the eddy viscosity [$cm^2 \cdot s^{-1}$]} \nonumber\\
Sc_T &= \mbox{the turbulent Schmidt number [-].} \nonumber
\end{align}

%The terms in Equations \eqref{eq:mg-n-diff} and \eqref{eq:mg-n-diff-Ci} are defined as follows:
%\begin{description}
%    \item $\nu_g$ is the average number of neutrons produced by fission of group $g$
%    \item $\phi_g$ is the scalar neutron flux of group $g$
%    \item $D_g$ is the group $g$ diffusion term
%    \item $G$ is the total number of energy groups
%    \item $\Sigma_{s, g' \rightarrow g}$ is the macroscopic scattering cross section from group $g'$ to group $g$
%    \item $\beta$ is the delayed neutron fraction
%    \item $\chi_{p, g}$ is the number of prompt neutrons born in group $g$
%    \item $\Sigma_f$ is the macroscopic fission cross section
%    \item $I$ is the number of delayed neutron precursor groups
%    \item $\chi_{d, g, i}$ is the number of delayed neutrons from precursor group $i$ born in group $g$
%    \item $\lambda_i$ is the decay constant for group $i$
%    \item $C_i$ is the concentration of group $i$
%    \item $\vec{u}$ is the velocity profile
%    \item $\Sigma_{a, g}$ is the macroscopic absorption cross section for group $g$
%    \item $\beta_i$ is the delayed neutron fraction for group $i$
%    \item $\nu_T$ is the eddy viscosity
%    \item $Sc_T$ is the turbulent Schmidt number
%\end{description}


\subsubsection{Modified Point Kinetics}

\Gls{PRK} are a simpler approach to modeling a transient, with the basic form given as:
\begin{equation}
    \frac{dn(t)}{dt} = \frac{\rho(t) - \beta_{eff}}{\Lambda} n(t) + \sum_{i=1}^I \lambda_i C_i(t),
    \label{eq:prk-n}
\end{equation}
and,
\begin{equation}
    \frac{dC_i(t)}{dt} = \frac{\beta_{eff, i}}{\Lambda} n(t) - \lambda_i C_i(t),
    \label{eq:prk-C}
\end{equation}
where
\begin{align}
n &= \mbox{ the relative neutron density [-]} \nonumber\\
\rho &= \mbox{ the reactivity [-]} \nonumber\\
\beta_{eff} &= \mbox{ the effective delayed neutron fraction [-].} \nonumber\\
\end{align}
The derivation of the \gls{PRK} equations can be found in the work by Henry, with a fluid-fueled derivation in the work by Lapenta et al. \cite{henry1958application, lapenta2001point}.
\Gls{MPK} are used for \glspl{MSR}, as standard \gls{PRK} do not properly handle the movement of the fuel salt.



%\begin{description}
%    \item $n$ is the relative neutron density
%    \item $\rho$ is the reactivity
%    \item $\beta_{eff}$ is the effective delayed neutron fraction
%    \item $\Lambda$ is the prompt neutron generation time
%    \item $\lambda_i$ is the group $i$ decay constant
%    \item $C_i$ is the concentration of group $i$
%    \item $\beta_{eff, i}$ is the effective delayed neutron fraction from group $i$
%\end{description}

\Gls{PRK} generally struggle to capture spatial effects, such as flow variation, which is important for \glspl{MSR}. 
\Gls{PRK} do, however, give a correct prediction for general evolution trends \cite{wooten2018review}.
Wooten et al. provide four different publications which use \gls{MPK} in \gls{MSR} analysis \cite{guo2013simulations, merle2015physical, shi2016development, zanetti2015extension}.

The work by Guo et al. analyzes \gls{ULOHS} and \gls{ULOF} transients in \gls{MOSART}.
The \gls{PRK} equations are modified into the \gls{DPK}, which model precursor concentrations in-core and ex-core, with the primary set of equations given as:
\begin{equation}
    \frac{dn(t)}{dt} = \frac{\rho - \tilde{\beta}_{eff}}{\Lambda} n(t) + \sum_{i=1}^{I} \lambda_i C_{i, c} (t),
    \label{eq:dpk-n}
\end{equation}
\begin{equation}
\begin{split}
    \frac{dC_{i, c}(t)}{dt} =& \frac{\beta_i}{\Lambda} n(t) - \left(\lambda_i + \frac{1}{\tau_{c}} \right) C_{i, c}(t)\\
    & + \frac{1}{\tau_{l}} \left(\frac{V_{l}}{V_{c}}\right) C_{i, l} (t - \tau_{l}),
    \label{eq:dpk-cin}
\end{split}
\end{equation}
\begin{equation}
    \frac{dC_{i, l}(t)}{dt} = \frac{1}{\tau_{c}} \left(\frac{V_{c}}{V_{l}}\right) C_{i, c} (t) - \left(\lambda_i + \frac{1}{\tau_{l}} \right) C_{i, l}(t).
    \label{eq:dpk-cex}
\end{equation}
These equations, shown in Equations \eqref{eq:dpk-n} through \eqref{eq:dpk-cex} (recreated from \cite{guo2013simulations, shi2016development}), are also used by Shi et al. in \gls{RELAP}.
Shi et al. simulate a load demand change, primary flow transient, secondary flow transient, and reactivity transient in the \gls{MSBR} \cite{shi2016development}.
The modified point kinetics approach enables modeling of ex-core decay of \glspl{DNP}, capturing effects during power operation and loss-of-flow conditions, but it does not accurately capture all oscillatory effects \cite{wooten2018review}.

Intermediate terms used in Equations \eqref{eq:dpk-n} through \eqref{eq:dpk-cex} are shown in Equations \eqref{eq:dpk-betaeff} through \eqref{eq:dpk-tauloop} as:
\begin{equation}
    \tilde{\beta}_{eff} = \beta_{static} - \beta_{loss},
    \label{eq:dpk-betaeff}
\end{equation}
\begin{equation}
    \beta_{loss} = \frac{\Lambda(t)}{n(t)} \sum_{i=1}^I \lambda_i C_{i, l},
    \label{eq:dpk-betaloss}
\end{equation}
\begin{equation}
    \tau_{c} = \frac{M_{c}}{\dot{M}(t)},
    \label{eq:dpk-taucore}
\end{equation}
\begin{equation}
    \tau_{l} = \frac{M_{l}}{\dot{M}(t)}.
    \label{eq:dpk-tauloop}
\end{equation}

The terms in Equations \eqref{eq:dpk-n} through \eqref{eq:dpk-tauloop} are defined as follows:
\begin{align}
V_{l} &= \mbox{ the volume of the loop [$cm^3$]} \nonumber \\ 
V_{c} &= \mbox{ the volume of the core [$cm^3$]} \nonumber \\ 
\beta_{static} &= \mbox{ the value of $\beta_{eff}$ without flow [-]} \nonumber \\ 
M_{c} &= \mbox{ the mass in the core [$kg$]} \nonumber \\ 
M_{l} &= \mbox{ the mass in the loop [$kg$]} \nonumber \\ 
\dot{M}(t) &= \mbox{ the mass flow rate [$kg \cdot s^{-1}$].} \nonumber
\end{align}


%\begin{description}
%    \item $V_{loop}$ is the volume of the loop
%    \item $V_{core}$ is the volume of the core
%    \item $\beta_{static}$ is the value of $\beta_{eff}$ without flow
%    \item $M_{core}$ is the mass in the core
%    \item $M_{loop}$ is the mass in the loop
%    \item $w(t)$ is the mass flow rate
%\end{description}

A similar method is also used by Suzuki and Shimazu for analyzing a reactivity initiated accident in the FUJI-12 reactor \cite{suzuki2008reactivity}.
That work uses a simpler form of the equations, without dividing the precursor concentrations into core and loop regions.

The work by Merle-Lucotte et al. uses the "IPK method," which may stand for "Improved Point-Kinetic method," though this is not explicitly stated \cite{merle2015physical}.
This method uses a fixed mesh for neutronics and a mobile mesh for fluid properties.
Merle-Lucotte et al. found that the \gls{PRK} method without any modifications was able to match the IPK method for general trends and end-states, but was unable to capture reactivity oscillations properly.
The IPK method was able to capture these oscillations, but assumes uniform fission and instantaneous heat transfer \cite{merle2015physical}.

The work by Zanetti et al. extends the TRACE code, a finite volume code used for \gls{LWR} transients \cite{zanetti2015extension}.
The \gls{DNP} concentration, dependent on both space and time, is discretized with an upwind scheme in space and a backward Euler method in time.
Adjoint-weighted quantities discretized with the trapezoidal rule are then passed to the \gls{PRK} equation.
The set of \gls{DNP} concentration and power, or neutron density, equations are resolved using LU factorization.

Another approach to modify the \gls{PRK} is to alter the kinetics parameters based on a phase-space integration weighted with the adjoint flux \cite{lapenta2000basic}.
Alternatively, effective flow-inclusive group yields can be calculated, as shown by Haubenreich and used in \gls{PRK} equations \cite{haubenreich1962prediction}.

\subsubsection{Quasi-Statics}

Quasi-static methods seek to maintain the accuracy of multi-group diffusion approaches while also having the cost savings of \gls{PRK}.
The general approach of these methods is to update \gls{PRK} parameters with occasional multi-group diffusion simulations to capture changes in flux and \gls{DNP} spatial distributions \cite{wooten2018review}.
The general formulation used for quasi-static methods is very similar among the various publications that use it \cite{dulla2004neutron, tano2023coupled}.
A complete derivation of the quasi-static method can be found in the work by Henry, who discusses the shape and amplitude flux functions, and the work by Devooght, who discusses the  generalized quasi-static method \cite{henry1958application, devooght1980quasistatic}.
There are also other studies which discuss the quasi-static method in detail \cite{becker1968generalized}.

Wooten et al. introduce several different publications which have used quasi-statics \cite{wooten2018review}.
One of these, from 2005 by Rineiski et al., covers the SIMMER code, which couples neutronics, thermal-hydraulics, and structure codes \cite{rineiski2005kinetics}.
The neutron flux is constrained in the work by a term containing the phase-space integration and the weighting function used.
The neutron flux is treated as:
\begin{equation}
    \Phi(\vec{x}, t) = \phi_A(t) \psi (\vec{x}, t),
    \label{eq:flux-treatment-rineiski}
\end{equation}
where
\begin{align}
    \Phi &= \text{the neutron flux [$n \cdot cm^{-2} \cdot s^{-1}$] } \nonumber \\
    \phi_A &= \text{the neutron flux amplitude function [n]} \nonumber \\
    \psi (\vec{x}, t) &= \text {the neutron flux shape function [$cm^{-2} \cdot s^{-1}$]. }\nonumber 
\end{align}

Rineiski et al. solve "shape" steps to handle the neutronics and \gls{DNP} spatial concentrations.
The equations used in the "shape" steps are recreated here in Equations \eqref{eq:rineiski-flux} and \eqref{eq:rineiski-big-C}.
These equations are solved by time-discretizing Equation \eqref{eq:rineiski-flux} into a set of steady-state-like equations, which are then treated using the DANTSYS module of SIMMER \cite{alcouffe1995dantsys}.
These primary equations are given as:
\begin{equation}
\begin{split}
    \frac{1}{\bar{\nu}} \frac{\partial \psi (\vec{x}, t)}{\partial t} + \frac{1}{\bar{\nu} \phi_A} \frac{d \phi_A}{dt} \psi (\vec{x}, t) + M(t) \psi (\vec{x}, t) = & \\
    (1 - \beta) \chi_p F(t) \psi (\vec{x}, t)
    + \frac{1}{\phi_A} \chi_d \sum_{i=1}^I \lambda_i C_i (\vec{x}, t) + \frac{1}{\phi_A} Q (\vec{x}, t) & ,
    \label{eq:rineiski-flux}
\end{split}
\end{equation}
and
\begin{equation}
    \frac{\partial C_i (\vec{x}, t)}{\partial t} = \phi_A \beta_i F \psi (\vec{x}, t) - \lambda_i C_i (\vec{x}, t) - \nabla ( \vec{u} (\vec{x}, t) C_i (\vec{x}, t)).
    \label{eq:rineiski-big-C}
\end{equation}

Within these "shape" steps are also sub-steps referred to as "reactivity" steps.
At each of these "reactivity" steps, the \gls{PRK} parameters and flux amplitude are calculated.
The equations used in the "reactivity" steps are recreated here as:
\begin{equation}
    \frac{d \phi_A}{dt} = \left( \frac{\rho - \beta_{eff}}{\Lambda} \right) \phi_A + q_{eff} + \sum_{i=1}^I \lambda_i \left( c_{s, i} + c_{m, i} \right),
    \label{eq:rineiski-n}
\end{equation}
\begin{equation}
    \frac{dc_{s, i}}{dt} = \frac{\beta_{eff}}{\Lambda} \phi_A - \lambda_i C_{s, i}.
    \label{eq:rineiski-cs}
\end{equation}
\begin{equation}
    \rho = \frac{\left\langle W(\vec{x}), \left(\chi F(t) - M(t)\right) \psi  (\vec{x}, t) \right\rangle}{\left\langle W(\vec{x}), \chi F(t) \psi  (\vec{x}, t) \right\rangle}.
    \label{eq:rineiski-rho}
\end{equation}
\begin{equation}
    c_{m, i} = \frac{\left\langle W(\vec{x}), \chi_d C_{m, i} \right\rangle}{\Lambda \left\langle W(\vec{x}), \chi F(t) \psi  (\vec{x}, t) \right\rangle}.
    \label{eq:rineiski-cm}
\end{equation}

Intermediate terms used in Equations \eqref{eq:rineiski-n} through \eqref{eq:rineiski-cm} are given as:
\begin{equation}
    \frac{\partial C_{m, i}(\vec{x}, t)}{\partial t} = - \nabla \left( \vec{u}(\vec{x}, t) C_{s, i} (\vec{x}, t) \right) - \lambda_i C_{m, i}(\vec{x}, t) - \nabla \left( \vec{u}(\vec{x}, t) C_{m, i}(\vec{x}, t)\right),
    \label{eq:rineiski-big-cm}
\end{equation}
and
\begin{equation}
    \frac{\partial C_{s, i}(\vec{x}, t)}{\partial t} = \phi_A(t) \beta_i F(t) \psi (\vec{x}, t) - \lambda_i C_{s, i}(\vec{x}, t).
    \label{eq:rineiski-big-cs}
\end{equation}
The $C_s(\vec{x}, t)$ term is computed at steady state by setting the left-hand side of both Equation \eqref{eq:rineiski-flux} and \eqref{eq:rineiski-big-C} to zero as well as the fluid velocity to zero at t=0.
This allows Equation \eqref{eq:rineiski-big-cm} to be solved.

The terms in Equations \eqref{eq:flux-treatment-rineiski} through \eqref{eq:rineiski-big-cs} are defined as follows:
\begin{align}
M(t) &= \mbox{ the neutron disappearance operator [$cm^{-1}$] } \nonumber \\
F(t) &= \mbox{ the neutron production operator [$cm^{-1}$]} \nonumber \\
\chi_p &= \mbox{ the prompt neutron energy spectra [-]} \nonumber \\
C_i(\vec{x}, t) &= \mbox{ the concentration of \gls{DNP} group \textit{i} [$atoms \cdot cm^{-3}$] } \nonumber \\
Q(\vec{x}, t) &= \mbox{ the external neutron source [$n \cdot cm^{-3} \cdot s^{-1}$] } \nonumber \\
\vec{u}(\vec{x}, t) &= \mbox{ the \gls{DNP} velocity [$cm \cdot s^{-1}$]} \nonumber \\
C_{s, i}(\vec{x}, t) &= \mbox{ the "standard" precursor concentration of group \textit{i} [$atoms \cdot cm^{-3}$] } \nonumber \\
C_{m, i}(\vec{x}, t) &= \mbox{ the "movable" precursor concentration of group \textit{i} [$atoms \cdot cm^{-3}$] } \nonumber \\
\rho &= \mbox{ the reactivity [-]} \nonumber \\
c_{s, i} &= \mbox{ the \gls{PRK} "standard" precursor concentration of group $i$ [$atoms$]} \nonumber \\
c_{m, i} &= \mbox{ the \gls{PRK} "movable" precursor concentration of group $i$ [$atoms$]} \nonumber \\
q_{eff} &= \mbox{ the \gls{PRK} external neutron source [$n \cdot s^{-1}$] } \nonumber \\
\Lambda &= \mbox{ the neutron generation time [$s$]} \nonumber \\
W(x) &= \mbox{ the weighting function (the default in SIMMER} \nonumber\\
&\;\;\;\;\; \mbox{ is an adjoint flux computed at steady-state conditions) [-].} \nonumber
\end{align}


%\begin{description}
%    \item $\Phi$ is the neutron flux
%    \item $N(t)$ is the neutron flux amplitude
%    \item $\psi(\vec{x}, t)$ is the neutron flux shape
%    \item $\nu$ is the average neutron speed
%    \item $M(t)$ is the neutron disappearance operator
%    \item $F(t)$ is the neutron production operator
%    \item $\beta$ is the delayed neutron fraction
%    \item $\beta_{eff}$ is the effective delayed neutron fraction
%    \item $\chi_p$ is the prompt neutron energy spectra
%    \item $\chi_d$ is the delayed neutron energy spectra
%    \item $\chi$ is the neutron energy spectra
%    \item $\lambda_i$ is the decay constant for \gls{DNP} group $i$
%    \item $C_i(\vec{x}, t)$ is the concentration of \gls{DNP} group $i$
%    \item $Q(\vec{x}, t)$ is the external neutron source
%    \item $u(\vec{x}, t)$ is the \gls{DNP} velocity
%    \item $C_{s, i}(\vec{x}, t)$ is the "standard" precursor concentration of group $i$
%    \item $C_{m, i}(\vec{x}, t)$ is the "movable" precursor concentration of group $i$
%    \item $\rho$ is the reactivity
%    \item $c_{s, i}$ is the \gls{PRK} "standard" precursor concentration of group $i$
%    \item $c_{m, i}$ is the \gls{PRK} "movable" precursor concentration of group $i$
%    \item $q_{eff}$ is the \gls{PRK} external neutron source 
%    \item $\Lambda$ is the prompt neutron generation time
%    \item $W(x)$ is the weighting function (by default in SIMMER is an adjoint flux computed at steady-state conditions)
%\end{description}



%\subsubsection{Alternative Methods}

%Although multi-group diffusion, \gls{PRK}, and quasi-statics are the most commonly used methods, there are other choices used to analyze transients.
%One of these is the \gls{TFM} method, which entails generating fission matrices for use in a \gls{CFD} code to model transients \cite{laureau2015coupled}.


